\subsection{Reproduction and the effect of the protocol}
Table~\ref{tab:protocol} and Fig.~\ref{fig:protocol} compare protocols for the model trained on tumour-free BraTS slices.
Under REFLECT's protocol our pooled Dice is 81.7 (80.8 with one step), against 79.6 reported; the mean per-slice Dice is 74.5.
The paper does not state which aggregate it reports, so the reproduction lies within this range.
Fixing the threshold on validation subjects instead of the test slices costs almost nothing (79.6 against a test-optimal 79.7).
Scoring all 20 central slices changes pooled Dice little, because large tumours dominate the pixel count, but the mean Dice over tumour slices falls to 64.1, and 7.7\% of tumour-free slices contain at least one false-positive pixel.

\begin{table}[t]
\centering\small
\caption{REFLECT on BraTS 2021 T2 (trained on tumour-free BraTS slices) under the original and strict protocols. Brackets: 95\% bootstrap intervals over test subjects.}
\label{tab:protocol}
\begin{tabular}{lll}
\toprule
Protocol & Dice & Notes \\
\midrule
Paper \citep{beizaee2025reflect}, REFLECT protocol & 79.6 & \\
Ours, REFLECT protocol, pooled & 81.7 & 1 step: 80.8 \\
Ours, REFLECT protocol, mean per slice & 74.5 & 1 step: 73.7 \\
Ours, strict, pooled & 79.6 [76.9, 81.9] & validation threshold \\
Ours, strict, mean over tumour slices & 64.1 [59.7, 68.3] & \\
\bottomrule
\end{tabular}
\end{table}

\begin{figure}[t]
\centering
\includegraphics[width=\linewidth]{protocol.pdf}
\caption{Left: Dice of the BraTS-trained and IXI-trained models under four protocols. Right: false-alarm rate on tumour-free slices, which the single-slice protocol cannot measure.}
\label{fig:protocol}
\end{figure}

\subsection{Healthy training data from other hospitals}
Table~\ref{tab:site} reports all models on the strict protocol.
Replacing tumour-free BraTS slices with IXI healthy volunteers lowers AUPRC from 0.865 to 0.828 (paired difference $-3.8$ points, 95\% interval $[-5.7, -2.0]$) and pooled Dice at each model's Dice-optimal threshold by 3.5 points $[2.1, 5.0]$.
The same thresholds, however, let false alarms rise from 7.7\% to 29.0\% of tumour-free slices ($+21.3$ $[9.3, 35.3]$).
Each model thus reaches its best Dice at a different false-alarm rate, and comparing Dice-optimal scores understates the shift.
Comparing the models where each fires on 10\% of tumour-free test slices, the Dice gap is 5.7 points $[2.7, 12.3]$; at 5\% it is 11.6 points $[3.1, 25.4]$.
Fig.~\ref{fig:tradeoff} shows the full trade-off: the curves nearly meet at high false-alarm rates and separate at the low rates a screening tool would operate at.
Thresholds chosen on validation for a 5\% false-alarm rate gave 7--17\% on test, so validation false-alarm budgets transfer only roughly between subject sets.

\begin{table}[t]
\centering\small
\setlength{\tabcolsep}{4pt}
\caption{Strict protocol, 5 correction steps. FA: tumour-free test slices with any positive pixel (\%). ``10\% test FA'': threshold at which 10\% of tumour-free test slices fire, chosen per model and per bootstrap resample. Brackets: 95\% bootstrap intervals over the 126 test subjects. Paired differences are given in the text.}
\label{tab:site}
\begin{tabular}{lccccc}
\toprule
& & \multicolumn{2}{c}{Dice-optimal validation threshold} & \multicolumn{2}{c}{10\% test FA} \\
\cmidrule(lr){3-4}\cmidrule(lr){5-6}
Healthy training data & AUPRC & pooled Dice & FA & pooled Dice & mean Dice \\
\midrule
BraTS (in-domain) & 0.865 & 79.6 [76.9, 81.9] & 7.7 [2.5, 15.1] & 79.6 [76.8, 81.7] & 64.3 [59.5, 68.5] \\
IXI & 0.828 & 76.1 [72.6, 79.0] & 29.0 [17.2, 43.5] & 73.9 [66.3, 77.9] & 57.0 [47.8, 62.8] \\
IXI, thick-slice aug. & 0.862 & 77.7 [74.5, 80.5] & 42.3 [29.2, 57.2] & 67.6 [54.5, 75.6] & 49.1 [35.9, 59.1] \\
IXI, +10 ep.\ IXI only & 0.828 & 75.8 [72.5, 78.7] & 31.1 [19.3, 45.8] & 72.5 [66.1, 76.9] & 55.3 [48.6, 61.8] \\
IXI, +10 ep.\ 5 local & 0.829 & 76.4 [73.0, 79.1] & 27.7 [16.5, 41.8] & 73.0 [66.8, 77.3] & 55.6 [48.5, 61.7] \\
IXI, +10 ep.\ 20 local & 0.833 & 76.9 [73.8, 79.5] & 23.7 [14.0, 35.8] & 75.9 [71.9, 79.3] & 59.0 [54.0, 64.4] \\
IXI, +10 ep.\ 50 local & 0.834 & 76.8 [73.8, 79.3] & 22.9 [13.1, 35.0] & 76.3 [72.3, 79.3] & 60.4 [54.8, 65.2] \\
\bottomrule
\end{tabular}
\end{table}

\begin{figure}[t]
\centering
\includegraphics[width=\linewidth]{tradeoff.pdf}
\caption{Pooled Dice (left) and mean Dice over tumour slices (right) against the false-alarm rate on tumour-free test slices, over all thresholds. Crosses: Dice-optimal validation threshold; dots: threshold for at most 5\% validation false alarms.}
\label{fig:tradeoff}
\end{figure}

\subsection{Where the false alarms are}
On tumour-free test slices, the IXI model's false-positive pixels are bright (mean normalised intensity 0.45 against 0.27 for the brain), and only about 10\% lie within 8 pixels of the brain border, for both models.
A subject's false-alarm rate does not correlate with a through-plane sharpness proxy (Spearman $\rho=0.19$, $p=0.24$, 41 subjects), and 21.5\% of slices fire even in subjects whose 20 central slices contain no tumour, so neither slice thickness nor proximity to the tumour explains them.
Ventricle size and central CSF fraction show no significant correlation either, although IXI subjects have larger central CSF than BraTS subjects ($p<10^{-19}$).
Visually (Fig.~\ref{fig:qual} and the repository's \texttt{fp\_slices.png}), the flagged regions are ventricles enlarged or displaced by mass effect, low frontal slices near the skull base, and distorted anatomy.
The model trained on tumour-free slices of the BraTS patients reconstructs all of these.
We read this as anatomy leaking from patients into the in-domain ``healthy'' training data: those slices come from tumour patients whose brains are deformed by mass effect from the tumour elsewhere in the volume, which healthy volunteers do not show.

\begin{figure}[t]
\centering
\includegraphics[width=\linewidth]{qualitative.pdf}
\caption{Test slices, healthy reconstructions and thresholded anomaly maps (threshold for at most 5\% validation false alarms; tumour outlined in cyan). Tumour rows are drawn at random from three size quantiles; the two tumour-free rows are the slices on which the IXI model fires most, chosen to show its failure mode.}
\label{fig:qual}
\end{figure}

\subsection{Remedies}
\paragraph{Intensity harmonisation.}
Matching each test subject's brain intensity histogram to the IXI training distribution before inference raises false alarms to 33.5\% and lowers pooled Dice to 75.0; the shift is not one of intensity distribution alone.

\paragraph{Thick-slice augmentation.}
Many BraTS T2 scans are thick-slice acquisitions resampled to 1\,mm.
We simulate this on IXI by averaging slabs of 1--6\,mm along a random axis and resampling.
AUPRC rises from 0.828 to 0.862 ($+3.4$ $[1.5, 5.5]$), nearly the in-domain 0.865, and Dice at the Dice-optimal threshold from 76.1 to 77.7.
The model also flags more healthy tissue, however: 42.3\% false alarms at that threshold, and at a threshold fixed for 5\% validation false alarms its pooled Dice is 64.4 against 72.0 for plain IXI training ($-7.6$ $[-9.2, -6.0]$) at similar test false-alarm rates (8.2\% and 7.4\%).
Threshold-free pixel metrics and Dice-optimal Dice can therefore favour a model that is worse at a clinically usable operating point.

\paragraph{Local fine-tuning.}
We continue training the IXI model for 10 epochs (learning rate $5\times10^{-5}$) on IXI plus the tumour-free slices of $N$ randomly chosen BraTS training subjects, oversampled ten times; $N=0$ is a control with the same extra training on IXI only.
Fig.~\ref{fig:local} shows paired differences from the IXI model.
With 20 and 50 local subjects, false alarms at the Dice-optimal threshold fall by 5.3 $[0.7, 10.1]$ and 6.1 $[1.6, 10.9]$ points, and pooled Dice at a 10\% test false-alarm rate rises by 2.0 $[0.6, 7.7]$ and 2.4 $[0.6, 7.9]$ points, recovering 35--42\% of the gap to the in-domain model.
The control changes neither (intervals include zero), so the gain comes from the local data, not from extra training.
Five local subjects are not enough.

\begin{figure}[t]
\centering
\includegraphics[width=0.85\linewidth]{local_finetuning.pdf}
\caption{Local fine-tuning: paired differences from the IXI model in false alarms at the Dice-optimal threshold (left) and in pooled Dice at a 10\% test false-alarm rate (right). Bars: 95\% paired bootstrap intervals. Dashed: the in-domain model.}
\label{fig:local}
\end{figure}
